\section{Related work}
\label{sec:related}

\subsection{Automatic sleep staging}
\label{sec:related-staging}

Single-channel EEG sleep staging is dominated by compact convolutional and
recurrent architectures. \citet{lawhern2018} introduced EEGNet, a
depthwise-separable convolutional network that remains a standard baseline;
\citet{supratak2017} and \citet{phan2019seqsleepnet} model the sequence of
epochs with recurrent layers over per-epoch convolutional features;
\citet{perslev2021} use a fully convolutional network over the whole night;
and \citet{phan2022sleepformer} replace recurrence with self-attention and
add an entropy-based uncertainty estimate. Recent work in this journal
continues the same line on the same corpus, with rhythm-aware contrastive
representation learning \citep{zhou2026bspc} and adaptive multi-scale
convolution with local attention \citep{fan2025bspc}, and statistical
modelling of the sleep EEG itself \citep{hatamoto2025}. End-to-end networks
can learn directly from raw single-channel EEG \citep{li2021endtoend}, and
hybrid convolutional--HMM designs recover temporal continuity for point
prediction \citep{yang2021hmm}; neither construction supplies a calibrated
set whose coverage is evaluated per patient. A broad evaluation of
contemporary staging methods also illustrates how strongly reported
performance depends on the dataset, preprocessing, and evaluation protocol
\citep{sarkar2022evaluation}. Reported accuracy on
Sleep-EDF is typically \SIrange{82}{87}{\percent} with Cohen's kappa around
$0.75$--$0.82$; N1 remains the hardest class, with F1 often below $0.5$,
reflecting the poor inter-scorer agreement of that stage
\citep{rosenberg2013}.

Evaluation in this literature is usually epoch-level accuracy under a
subject-wise split, sometimes with a separate held-out cohort. Uncertainty,
when reported, is reported as a calibration error
\citep{guo2017,phan2022sleepformer} rather than as a coverage guarantee, and
almost always in aggregate over epochs.

\subsection{Uncertainty quantification for sleep staging}
\label{sec:related-uq}

Monte-Carlo dropout, deep ensembles, evidential deep learning and Bayesian
last layers all produce a per-epoch confidence, and \citet{phan2022sleepformer}
use predictive entropy to flag epochs for review. These scores rank epochs by
difficulty but carry no statement that the true stage is inside the returned
set with a given probability. Selective prediction for sleep staging has been
studied by \citet{cover2026}, who learn an error-ranking score and accept a
fixed fraction of epochs subject to an empirical risk target, deferring the
remainder to a larger model; this is complementary to our work, since
selective prediction returns a point label plus an abstention whereas
conformal prediction returns a set with a distribution-free coverage
guarantee, but the risk control is empirical rather than finite-sample.
Conformal prediction has also been embedded as a component of a
multiple-instance sleep pipeline \citep{conmil2025}, without a study of
coverage under subject shift. Uncertainty-guided pipelines such as U-PASS
use uncertainty operationally during sleep-model development
\citep{heremans2024upass}; the present work instead evaluates the realised
coverage and workload of the deployed set predictor. Recent work on sleep
staging has emphasised the distinction between aleatoric uncertainty,
reflecting genuine ambiguity among scorers or signals, and epistemic
uncertainty, reflecting limited model knowledge \citep{vangorp2022certainty}.
Conformal sets address the second object directly by controlling the error
rate of a predictive set, while their purity and workload reveal the first.

\paragraph{The evaluation gap}
Expected calibration error, maximum calibration error and the Brier score
measure the agreement between confidence and accuracy in aggregate. They say
nothing about whether a \emph{set} contains the truth, and they are computed
on the same pooled sample that hides subject heterogeneity. In our
experiments the EEGNet backbone, with an ECE of $0.030$, still leaves its
worst patient $12.8$ points below nominal coverage, which is the practical
failure mode that calibration metrics do not detect.

\subsection{Conformal prediction}
\label{sec:related-cp}

Split conformal prediction \citep{vovk2005,lei2018,angelopoulos2023gentle}
gives finite-sample marginal coverage under exchangeability. Adaptive
prediction sets (APS) and their regularised variant (RAPS)
\citep{romano2020,angelopoulos2021raps} adapt set size to the difficulty of
each example; Mondrian and class-conditional CP \citep{vovk2012} give
per-class coverage, which on our cohort shrinks sets rather than enlarging
them but leaves a residual patient gap because patients differ in ways that
are not their stage mix; weighted CP
\citep{tibshirani2019} handles a known covariate shift. None of these
provides a per-patient guarantee when the exchangeable unit is the patient
rather than the epoch, and \citet{barber2023} show that without
exchangeability only a weaker, non-vanishing bound is available.

Adaptive conformal inference \citep{gibbs2021} replaces exchangeability with a
long-run coverage guarantee along a stream, updating the target miscoverage
after each observed label; \citet{angelopoulos2024} analyse the step-size schedule
and show that decaying steps with $\sum_t\gamma_t=\infty$ and
$\sum_t\gamma_t^2<\infty$ give the strongest guarantee. Standard ACI runs a
single stream over all observations, so its guarantee is again pooled; we
instead run one stream per patient, which is the natural unit of clinical
deployment and which, to our knowledge, has not been evaluated for
five-class sleep staging.

\paragraph{Structured and non-exchangeable data}
\citet{nettasinghe2023} extend CP to hidden Markov models by partitioning the
process into exchangeable blocks via de Finetti's theorem for Markov chains,
producing a confidence set for a sequence through the filtering recursion.
Our block-level construction shares the motivation --- blocks, not epochs,
are the exchangeable unit --- but differs in what is produced: they construct
a set of sequences, whereas we project a scalar block-level conformal set onto
each position with a forward--backward dynamic program, which gives
simultaneous per-position coverage at $O(LK^2)$. We show in
\Cref{sec:results-temporal} that this projection, while valid, is too loose
to be useful, which is itself informative for anyone considering the
construction.

\paragraph{Time-series conformal prediction}
A parallel literature targets real-valued forecasting under temporal
dependence. EnbPI \citep{xu2021} bootstraps the residuals of a sequential
predictor and calibrates a quantile of their running maximum, so the interval
adapts to serial correlation without exchangeability; CopulaCPTS
\citep{stankeviciute2021} models the joint distribution of multi-step
residuals with a copula so that coverage holds for the whole forecast path;
and \citet{zaffran2022} update the threshold online for non-stationary
series. Our target is a categorical sequence with strong regime structure,
and the temporal information enters through the HMM posterior rather than
through a residual process. The path-level coverage those methods target is
what our SeqCP-Block constructor formalises, and
\Cref{sec:results-temporal} reports that on sleep blocks it is vacuous.

\paragraph{Healthcare and EEG}
\citet{chatterjee2026} show that split CP on EEG event classification loses
over 20 coverage points under patient shift and that per-patient
recalibration recovers most of it; this is the closest prior result to our
adaptive component, and we treat it as such. Our contribution there is not
the observation that personalisation helps but the deployment protocol: how
much to personalise, chosen by leave-one-subject-out simulation on the
calibration cohort rather than on the test patients; a streaming variant that
needs only one label at a time; and an evaluation that reads coverage against
set size for a five-class, strongly imbalanced task under subject-wise folds.
